\documentclass[11pt,letterpaper]{article}
\usepackage[margin=1in]{geometry}
\usepackage[hidelinks]{hyperref}
\pagestyle{empty}
\setlength{\parindent}{0pt}
\setlength{\parskip}{0.8em}
\begin{document}
\textbf{To the Editors of PLOS ONE}

Dear Editors,

Please consider our Research Article, \emph{Certifying best-of-$N$ verifier comparisons with fallible answer judgments}. A mathematical verifier chooses one solution from several candidates. When two verifiers favor different answers, the outcome of a benchmark comparison can turn on a small number of grades. Which grades deserve an expensive second check?

Earlier work by Tan and Clarke bounds ranking differences with incomplete judgments. Schnabel and colleagues use differences in selection probabilities to guide sampling. We study what happens when answers already have preliminary grades and some may be wrong. Given an upper bound on their number, we show how to allocate ordinary and reliable checks to minimize the worst remaining uncertainty. Once the preliminary grades are visible, the direction in which an error would move the comparison determines which checks are most valuable.

We tested the practical premise on public mathematical reasoning benchmarks. In one full-pool comparison, assuming at most five grading errors excluded the archived gap of 5.8 percentage points under both provisional grading rules. After estimating the error ceiling from a random sample, all six primary intervals included zero. Direct probability sampling at the same check budget yielded narrower median intervals in every setting. Readers can use these results to decide when they have the evidence to report a worst-case guarantee and when a statistical comparison is more informative.

We suggest Issa Ali Atoum as Academic Editor given his work on machine learning evaluation and reproducible research artifacts. The accompanying files provide the scored-data sources, protocol, code and numerical results needed to check the claims.

% Insert confirmed prior PLOS correspondence and reviewer exclusions here.
Sincerely,\\
Hanqing Yang\\
Corresponding author\\
Independent Researcher, Suzhou, China\\
\href{mailto:hanqing.yang189@gmail.com}{hanqing.yang189@gmail.com}
\end{document}
